% ====================================================================
% Mapa tematico de literatura y brecha de investigacion
% ====================================================================
\begin{figure}[H]
\centering
\resizebox{0.96\textwidth}{!}{%
\begin{tikzpicture}[
  font=\sffamily\small,
  panel/.style={draw=VIUDark!65,fill=VIULight,rounded corners=5pt,
    align=center,text width=4.15cm,minimum height=2.85cm,inner sep=8pt},
  gap/.style={draw=VIUOrange,fill=VIUOrange!10,rounded corners=5pt,
    align=center,text width=4.75cm,minimum height=2.65cm,inner sep=8pt,
    font=\sffamily\small\bfseries},
  flow/.style={-{Latex[length=2.5mm]},line width=1.0pt,draw=VIUOrange!90!black}
]
\node[panel] (lit) at (0,0) {
  \textbf{Familias parciales}\\[3pt]
  AMR en almacenes\\
  MRTA y subastas\\
  Juegos poblacionales\\
  Transporte cooperativo
};

\node[gap] (g) at (5.45,0) {
  Brecha del TFM\\[2pt]
  una arquitectura auditable conecta\\
  qu\'orum, capacidad,\\
  wrench y movimiento
};

\node[panel] (close) at (10.9,0) {
  \textbf{Interfaz evaluada}\\[3pt]
  A0--FULL integrado\\
  juegos y cierres separados\\
  localidad declarada\\
  evidencia SP1--SP8
};

\draw[flow] (lit.east) -- (g.west);
\draw[flow] (g.east) -- (close.west);
\end{tikzpicture}}
\caption[Mapa de literatura y brecha de investigación]{Mapa temático de literatura y brecha de investigación. El TFM no propone una única ley que resuelva toda la cadena: hace auditables las interfaces entre asignación, juegos, cierre entero, factibilidad física y ejecución, y declara qué operaciones conservan información global.}
\viuownsource
\label{fig:gap_literatura}
\end{figure}
